% =====================================================================
% backmatter/anexos.tex
% PT: Anexos. Ao contrário dos apêndices, os anexos contêm documentos que
%     NÃO foram produzidos pelo aluno (regulamentos, fichas técnicas,
%     autorizações). Se não tiver anexos, apague a linha
%     % =====================================================================
% backmatter/anexos.tex
% PT: Anexos. Ao contrário dos apêndices, os anexos contêm documentos que
%     NÃO foram produzidos pelo aluno (regulamentos, fichas técnicas,
%     autorizações). Se não tiver anexos, apague a linha
%     % =====================================================================
% backmatter/anexos.tex
% PT: Anexos. Ao contrário dos apêndices, os anexos contêm documentos que
%     NÃO foram produzidos pelo aluno (regulamentos, fichas técnicas,
%     autorizações). Se não tiver anexos, apague a linha
%     % =====================================================================
% backmatter/anexos.tex
% PT: Anexos. Ao contrário dos apêndices, os anexos contêm documentos que
%     NÃO foram produzidos pelo aluno (regulamentos, fichas técnicas,
%     autorizações). Se não tiver anexos, apague a linha
%     \input{backmatter/anexos} em main.tex.
% EN: Annexes. Unlike appendices, annexes hold documents that were NOT
%     produced by the student (regulations, data sheets, authorisations).
%     If you have no annexes, delete the line \input{backmatter/anexos}
%     in main.tex.
% =====================================================================
\chapter{Annexes}
\label{app:annexes}

% PT: Exemplo em inglês
% EN: English example
Annexes are documents written by others and included for reference. A PDF
document can be inserted page by page with \texttt{\textbackslash
includepdf} from the \texttt{pdfpages} package, as in the commented example
in this file.

% PT: Exemplo em português
% EN: Portuguese example
\begin{otherlanguage}{portuguese}
Exemplo em português: os anexos são documentos de terceiros, incluídos para
consulta. Para inserir um PDF, coloque-o na pasta \texttt{figures/} (ou numa
pasta \texttt{annexes/}) e descomente a linha \texttt{\textbackslash
includepdf} deste ficheiro.
\end{otherlanguage}

% PT: Inserir todas as páginas de um PDF, à escala da página.
%     pages=- : todas as páginas; pages={1-3} : só as páginas 1 a 3.
% EN: Insert every page of a PDF, scaled to the page.
%     pages=- : all pages; pages={1-3} : only pages 1 to 3.
% \includepdf[pages=-, scale=0.9, pagecommand={}]{annexes/document.pdf}
 em main.tex.
% EN: Annexes. Unlike appendices, annexes hold documents that were NOT
%     produced by the student (regulations, data sheets, authorisations).
%     If you have no annexes, delete the line % =====================================================================
% backmatter/anexos.tex
% PT: Anexos. Ao contrário dos apêndices, os anexos contêm documentos que
%     NÃO foram produzidos pelo aluno (regulamentos, fichas técnicas,
%     autorizações). Se não tiver anexos, apague a linha
%     \input{backmatter/anexos} em main.tex.
% EN: Annexes. Unlike appendices, annexes hold documents that were NOT
%     produced by the student (regulations, data sheets, authorisations).
%     If you have no annexes, delete the line \input{backmatter/anexos}
%     in main.tex.
% =====================================================================
\chapter{Annexes}
\label{app:annexes}

% PT: Exemplo em inglês
% EN: English example
Annexes are documents written by others and included for reference. A PDF
document can be inserted page by page with \texttt{\textbackslash
includepdf} from the \texttt{pdfpages} package, as in the commented example
in this file.

% PT: Exemplo em português
% EN: Portuguese example
\begin{otherlanguage}{portuguese}
Exemplo em português: os anexos são documentos de terceiros, incluídos para
consulta. Para inserir um PDF, coloque-o na pasta \texttt{figures/} (ou numa
pasta \texttt{annexes/}) e descomente a linha \texttt{\textbackslash
includepdf} deste ficheiro.
\end{otherlanguage}

% PT: Inserir todas as páginas de um PDF, à escala da página.
%     pages=- : todas as páginas; pages={1-3} : só as páginas 1 a 3.
% EN: Insert every page of a PDF, scaled to the page.
%     pages=- : all pages; pages={1-3} : only pages 1 to 3.
% \includepdf[pages=-, scale=0.9, pagecommand={}]{annexes/document.pdf}

%     in main.tex.
% =====================================================================
\chapter{Annexes}
\label{app:annexes}

% PT: Exemplo em inglês
% EN: English example
Annexes are documents written by others and included for reference. A PDF
document can be inserted page by page with \texttt{\textbackslash
includepdf} from the \texttt{pdfpages} package, as in the commented example
in this file.

% PT: Exemplo em português
% EN: Portuguese example
\begin{otherlanguage}{portuguese}
Exemplo em português: os anexos são documentos de terceiros, incluídos para
consulta. Para inserir um PDF, coloque-o na pasta \texttt{figures/} (ou numa
pasta \texttt{annexes/}) e descomente a linha \texttt{\textbackslash
includepdf} deste ficheiro.
\end{otherlanguage}

% PT: Inserir todas as páginas de um PDF, à escala da página.
%     pages=- : todas as páginas; pages={1-3} : só as páginas 1 a 3.
% EN: Insert every page of a PDF, scaled to the page.
%     pages=- : all pages; pages={1-3} : only pages 1 to 3.
% \includepdf[pages=-, scale=0.9, pagecommand={}]{annexes/document.pdf}
 em main.tex.
% EN: Annexes. Unlike appendices, annexes hold documents that were NOT
%     produced by the student (regulations, data sheets, authorisations).
%     If you have no annexes, delete the line % =====================================================================
% backmatter/anexos.tex
% PT: Anexos. Ao contrário dos apêndices, os anexos contêm documentos que
%     NÃO foram produzidos pelo aluno (regulamentos, fichas técnicas,
%     autorizações). Se não tiver anexos, apague a linha
%     % =====================================================================
% backmatter/anexos.tex
% PT: Anexos. Ao contrário dos apêndices, os anexos contêm documentos que
%     NÃO foram produzidos pelo aluno (regulamentos, fichas técnicas,
%     autorizações). Se não tiver anexos, apague a linha
%     \input{backmatter/anexos} em main.tex.
% EN: Annexes. Unlike appendices, annexes hold documents that were NOT
%     produced by the student (regulations, data sheets, authorisations).
%     If you have no annexes, delete the line \input{backmatter/anexos}
%     in main.tex.
% =====================================================================
\chapter{Annexes}
\label{app:annexes}

% PT: Exemplo em inglês
% EN: English example
Annexes are documents written by others and included for reference. A PDF
document can be inserted page by page with \texttt{\textbackslash
includepdf} from the \texttt{pdfpages} package, as in the commented example
in this file.

% PT: Exemplo em português
% EN: Portuguese example
\begin{otherlanguage}{portuguese}
Exemplo em português: os anexos são documentos de terceiros, incluídos para
consulta. Para inserir um PDF, coloque-o na pasta \texttt{figures/} (ou numa
pasta \texttt{annexes/}) e descomente a linha \texttt{\textbackslash
includepdf} deste ficheiro.
\end{otherlanguage}

% PT: Inserir todas as páginas de um PDF, à escala da página.
%     pages=- : todas as páginas; pages={1-3} : só as páginas 1 a 3.
% EN: Insert every page of a PDF, scaled to the page.
%     pages=- : all pages; pages={1-3} : only pages 1 to 3.
% \includepdf[pages=-, scale=0.9, pagecommand={}]{annexes/document.pdf}
 em main.tex.
% EN: Annexes. Unlike appendices, annexes hold documents that were NOT
%     produced by the student (regulations, data sheets, authorisations).
%     If you have no annexes, delete the line % =====================================================================
% backmatter/anexos.tex
% PT: Anexos. Ao contrário dos apêndices, os anexos contêm documentos que
%     NÃO foram produzidos pelo aluno (regulamentos, fichas técnicas,
%     autorizações). Se não tiver anexos, apague a linha
%     \input{backmatter/anexos} em main.tex.
% EN: Annexes. Unlike appendices, annexes hold documents that were NOT
%     produced by the student (regulations, data sheets, authorisations).
%     If you have no annexes, delete the line \input{backmatter/anexos}
%     in main.tex.
% =====================================================================
\chapter{Annexes}
\label{app:annexes}

% PT: Exemplo em inglês
% EN: English example
Annexes are documents written by others and included for reference. A PDF
document can be inserted page by page with \texttt{\textbackslash
includepdf} from the \texttt{pdfpages} package, as in the commented example
in this file.

% PT: Exemplo em português
% EN: Portuguese example
\begin{otherlanguage}{portuguese}
Exemplo em português: os anexos são documentos de terceiros, incluídos para
consulta. Para inserir um PDF, coloque-o na pasta \texttt{figures/} (ou numa
pasta \texttt{annexes/}) e descomente a linha \texttt{\textbackslash
includepdf} deste ficheiro.
\end{otherlanguage}

% PT: Inserir todas as páginas de um PDF, à escala da página.
%     pages=- : todas as páginas; pages={1-3} : só as páginas 1 a 3.
% EN: Insert every page of a PDF, scaled to the page.
%     pages=- : all pages; pages={1-3} : only pages 1 to 3.
% \includepdf[pages=-, scale=0.9, pagecommand={}]{annexes/document.pdf}

%     in main.tex.
% =====================================================================
\chapter{Annexes}
\label{app:annexes}

% PT: Exemplo em inglês
% EN: English example
Annexes are documents written by others and included for reference. A PDF
document can be inserted page by page with \texttt{\textbackslash
includepdf} from the \texttt{pdfpages} package, as in the commented example
in this file.

% PT: Exemplo em português
% EN: Portuguese example
\begin{otherlanguage}{portuguese}
Exemplo em português: os anexos são documentos de terceiros, incluídos para
consulta. Para inserir um PDF, coloque-o na pasta \texttt{figures/} (ou numa
pasta \texttt{annexes/}) e descomente a linha \texttt{\textbackslash
includepdf} deste ficheiro.
\end{otherlanguage}

% PT: Inserir todas as páginas de um PDF, à escala da página.
%     pages=- : todas as páginas; pages={1-3} : só as páginas 1 a 3.
% EN: Insert every page of a PDF, scaled to the page.
%     pages=- : all pages; pages={1-3} : only pages 1 to 3.
% \includepdf[pages=-, scale=0.9, pagecommand={}]{annexes/document.pdf}

%     in main.tex.
% =====================================================================
\chapter{Annexes}
\label{app:annexes}

% PT: Exemplo em inglês
% EN: English example
Annexes are documents written by others and included for reference. A PDF
document can be inserted page by page with \texttt{\textbackslash
includepdf} from the \texttt{pdfpages} package, as in the commented example
in this file.

% PT: Exemplo em português
% EN: Portuguese example
\begin{otherlanguage}{portuguese}
Exemplo em português: os anexos são documentos de terceiros, incluídos para
consulta. Para inserir um PDF, coloque-o na pasta \texttt{figures/} (ou numa
pasta \texttt{annexes/}) e descomente a linha \texttt{\textbackslash
includepdf} deste ficheiro.
\end{otherlanguage}

% PT: Inserir todas as páginas de um PDF, à escala da página.
%     pages=- : todas as páginas; pages={1-3} : só as páginas 1 a 3.
% EN: Insert every page of a PDF, scaled to the page.
%     pages=- : all pages; pages={1-3} : only pages 1 to 3.
% \includepdf[pages=-, scale=0.9, pagecommand={}]{annexes/document.pdf}
 em main.tex.
% EN: Annexes. Unlike appendices, annexes hold documents that were NOT
%     produced by the student (regulations, data sheets, authorisations).
%     If you have no annexes, delete the line % =====================================================================
% backmatter/anexos.tex
% PT: Anexos. Ao contrário dos apêndices, os anexos contêm documentos que
%     NÃO foram produzidos pelo aluno (regulamentos, fichas técnicas,
%     autorizações). Se não tiver anexos, apague a linha
%     % =====================================================================
% backmatter/anexos.tex
% PT: Anexos. Ao contrário dos apêndices, os anexos contêm documentos que
%     NÃO foram produzidos pelo aluno (regulamentos, fichas técnicas,
%     autorizações). Se não tiver anexos, apague a linha
%     % =====================================================================
% backmatter/anexos.tex
% PT: Anexos. Ao contrário dos apêndices, os anexos contêm documentos que
%     NÃO foram produzidos pelo aluno (regulamentos, fichas técnicas,
%     autorizações). Se não tiver anexos, apague a linha
%     \input{backmatter/anexos} em main.tex.
% EN: Annexes. Unlike appendices, annexes hold documents that were NOT
%     produced by the student (regulations, data sheets, authorisations).
%     If you have no annexes, delete the line \input{backmatter/anexos}
%     in main.tex.
% =====================================================================
\chapter{Annexes}
\label{app:annexes}

% PT: Exemplo em inglês
% EN: English example
Annexes are documents written by others and included for reference. A PDF
document can be inserted page by page with \texttt{\textbackslash
includepdf} from the \texttt{pdfpages} package, as in the commented example
in this file.

% PT: Exemplo em português
% EN: Portuguese example
\begin{otherlanguage}{portuguese}
Exemplo em português: os anexos são documentos de terceiros, incluídos para
consulta. Para inserir um PDF, coloque-o na pasta \texttt{figures/} (ou numa
pasta \texttt{annexes/}) e descomente a linha \texttt{\textbackslash
includepdf} deste ficheiro.
\end{otherlanguage}

% PT: Inserir todas as páginas de um PDF, à escala da página.
%     pages=- : todas as páginas; pages={1-3} : só as páginas 1 a 3.
% EN: Insert every page of a PDF, scaled to the page.
%     pages=- : all pages; pages={1-3} : only pages 1 to 3.
% \includepdf[pages=-, scale=0.9, pagecommand={}]{annexes/document.pdf}
 em main.tex.
% EN: Annexes. Unlike appendices, annexes hold documents that were NOT
%     produced by the student (regulations, data sheets, authorisations).
%     If you have no annexes, delete the line % =====================================================================
% backmatter/anexos.tex
% PT: Anexos. Ao contrário dos apêndices, os anexos contêm documentos que
%     NÃO foram produzidos pelo aluno (regulamentos, fichas técnicas,
%     autorizações). Se não tiver anexos, apague a linha
%     \input{backmatter/anexos} em main.tex.
% EN: Annexes. Unlike appendices, annexes hold documents that were NOT
%     produced by the student (regulations, data sheets, authorisations).
%     If you have no annexes, delete the line \input{backmatter/anexos}
%     in main.tex.
% =====================================================================
\chapter{Annexes}
\label{app:annexes}

% PT: Exemplo em inglês
% EN: English example
Annexes are documents written by others and included for reference. A PDF
document can be inserted page by page with \texttt{\textbackslash
includepdf} from the \texttt{pdfpages} package, as in the commented example
in this file.

% PT: Exemplo em português
% EN: Portuguese example
\begin{otherlanguage}{portuguese}
Exemplo em português: os anexos são documentos de terceiros, incluídos para
consulta. Para inserir um PDF, coloque-o na pasta \texttt{figures/} (ou numa
pasta \texttt{annexes/}) e descomente a linha \texttt{\textbackslash
includepdf} deste ficheiro.
\end{otherlanguage}

% PT: Inserir todas as páginas de um PDF, à escala da página.
%     pages=- : todas as páginas; pages={1-3} : só as páginas 1 a 3.
% EN: Insert every page of a PDF, scaled to the page.
%     pages=- : all pages; pages={1-3} : only pages 1 to 3.
% \includepdf[pages=-, scale=0.9, pagecommand={}]{annexes/document.pdf}

%     in main.tex.
% =====================================================================
\chapter{Annexes}
\label{app:annexes}

% PT: Exemplo em inglês
% EN: English example
Annexes are documents written by others and included for reference. A PDF
document can be inserted page by page with \texttt{\textbackslash
includepdf} from the \texttt{pdfpages} package, as in the commented example
in this file.

% PT: Exemplo em português
% EN: Portuguese example
\begin{otherlanguage}{portuguese}
Exemplo em português: os anexos são documentos de terceiros, incluídos para
consulta. Para inserir um PDF, coloque-o na pasta \texttt{figures/} (ou numa
pasta \texttt{annexes/}) e descomente a linha \texttt{\textbackslash
includepdf} deste ficheiro.
\end{otherlanguage}

% PT: Inserir todas as páginas de um PDF, à escala da página.
%     pages=- : todas as páginas; pages={1-3} : só as páginas 1 a 3.
% EN: Insert every page of a PDF, scaled to the page.
%     pages=- : all pages; pages={1-3} : only pages 1 to 3.
% \includepdf[pages=-, scale=0.9, pagecommand={}]{annexes/document.pdf}
 em main.tex.
% EN: Annexes. Unlike appendices, annexes hold documents that were NOT
%     produced by the student (regulations, data sheets, authorisations).
%     If you have no annexes, delete the line % =====================================================================
% backmatter/anexos.tex
% PT: Anexos. Ao contrário dos apêndices, os anexos contêm documentos que
%     NÃO foram produzidos pelo aluno (regulamentos, fichas técnicas,
%     autorizações). Se não tiver anexos, apague a linha
%     % =====================================================================
% backmatter/anexos.tex
% PT: Anexos. Ao contrário dos apêndices, os anexos contêm documentos que
%     NÃO foram produzidos pelo aluno (regulamentos, fichas técnicas,
%     autorizações). Se não tiver anexos, apague a linha
%     \input{backmatter/anexos} em main.tex.
% EN: Annexes. Unlike appendices, annexes hold documents that were NOT
%     produced by the student (regulations, data sheets, authorisations).
%     If you have no annexes, delete the line \input{backmatter/anexos}
%     in main.tex.
% =====================================================================
\chapter{Annexes}
\label{app:annexes}

% PT: Exemplo em inglês
% EN: English example
Annexes are documents written by others and included for reference. A PDF
document can be inserted page by page with \texttt{\textbackslash
includepdf} from the \texttt{pdfpages} package, as in the commented example
in this file.

% PT: Exemplo em português
% EN: Portuguese example
\begin{otherlanguage}{portuguese}
Exemplo em português: os anexos são documentos de terceiros, incluídos para
consulta. Para inserir um PDF, coloque-o na pasta \texttt{figures/} (ou numa
pasta \texttt{annexes/}) e descomente a linha \texttt{\textbackslash
includepdf} deste ficheiro.
\end{otherlanguage}

% PT: Inserir todas as páginas de um PDF, à escala da página.
%     pages=- : todas as páginas; pages={1-3} : só as páginas 1 a 3.
% EN: Insert every page of a PDF, scaled to the page.
%     pages=- : all pages; pages={1-3} : only pages 1 to 3.
% \includepdf[pages=-, scale=0.9, pagecommand={}]{annexes/document.pdf}
 em main.tex.
% EN: Annexes. Unlike appendices, annexes hold documents that were NOT
%     produced by the student (regulations, data sheets, authorisations).
%     If you have no annexes, delete the line % =====================================================================
% backmatter/anexos.tex
% PT: Anexos. Ao contrário dos apêndices, os anexos contêm documentos que
%     NÃO foram produzidos pelo aluno (regulamentos, fichas técnicas,
%     autorizações). Se não tiver anexos, apague a linha
%     \input{backmatter/anexos} em main.tex.
% EN: Annexes. Unlike appendices, annexes hold documents that were NOT
%     produced by the student (regulations, data sheets, authorisations).
%     If you have no annexes, delete the line \input{backmatter/anexos}
%     in main.tex.
% =====================================================================
\chapter{Annexes}
\label{app:annexes}

% PT: Exemplo em inglês
% EN: English example
Annexes are documents written by others and included for reference. A PDF
document can be inserted page by page with \texttt{\textbackslash
includepdf} from the \texttt{pdfpages} package, as in the commented example
in this file.

% PT: Exemplo em português
% EN: Portuguese example
\begin{otherlanguage}{portuguese}
Exemplo em português: os anexos são documentos de terceiros, incluídos para
consulta. Para inserir um PDF, coloque-o na pasta \texttt{figures/} (ou numa
pasta \texttt{annexes/}) e descomente a linha \texttt{\textbackslash
includepdf} deste ficheiro.
\end{otherlanguage}

% PT: Inserir todas as páginas de um PDF, à escala da página.
%     pages=- : todas as páginas; pages={1-3} : só as páginas 1 a 3.
% EN: Insert every page of a PDF, scaled to the page.
%     pages=- : all pages; pages={1-3} : only pages 1 to 3.
% \includepdf[pages=-, scale=0.9, pagecommand={}]{annexes/document.pdf}

%     in main.tex.
% =====================================================================
\chapter{Annexes}
\label{app:annexes}

% PT: Exemplo em inglês
% EN: English example
Annexes are documents written by others and included for reference. A PDF
document can be inserted page by page with \texttt{\textbackslash
includepdf} from the \texttt{pdfpages} package, as in the commented example
in this file.

% PT: Exemplo em português
% EN: Portuguese example
\begin{otherlanguage}{portuguese}
Exemplo em português: os anexos são documentos de terceiros, incluídos para
consulta. Para inserir um PDF, coloque-o na pasta \texttt{figures/} (ou numa
pasta \texttt{annexes/}) e descomente a linha \texttt{\textbackslash
includepdf} deste ficheiro.
\end{otherlanguage}

% PT: Inserir todas as páginas de um PDF, à escala da página.
%     pages=- : todas as páginas; pages={1-3} : só as páginas 1 a 3.
% EN: Insert every page of a PDF, scaled to the page.
%     pages=- : all pages; pages={1-3} : only pages 1 to 3.
% \includepdf[pages=-, scale=0.9, pagecommand={}]{annexes/document.pdf}

%     in main.tex.
% =====================================================================
\chapter{Annexes}
\label{app:annexes}

% PT: Exemplo em inglês
% EN: English example
Annexes are documents written by others and included for reference. A PDF
document can be inserted page by page with \texttt{\textbackslash
includepdf} from the \texttt{pdfpages} package, as in the commented example
in this file.

% PT: Exemplo em português
% EN: Portuguese example
\begin{otherlanguage}{portuguese}
Exemplo em português: os anexos são documentos de terceiros, incluídos para
consulta. Para inserir um PDF, coloque-o na pasta \texttt{figures/} (ou numa
pasta \texttt{annexes/}) e descomente a linha \texttt{\textbackslash
includepdf} deste ficheiro.
\end{otherlanguage}

% PT: Inserir todas as páginas de um PDF, à escala da página.
%     pages=- : todas as páginas; pages={1-3} : só as páginas 1 a 3.
% EN: Insert every page of a PDF, scaled to the page.
%     pages=- : all pages; pages={1-3} : only pages 1 to 3.
% \includepdf[pages=-, scale=0.9, pagecommand={}]{annexes/document.pdf}

%     in main.tex.
% =====================================================================
\chapter{Annexes}
\label{app:annexes}

% PT: Exemplo em inglês
% EN: English example
Annexes are documents written by others and included for reference. A PDF
document can be inserted page by page with \texttt{\textbackslash
includepdf} from the \texttt{pdfpages} package, as in the commented example
in this file.

% PT: Exemplo em português
% EN: Portuguese example
\begin{otherlanguage}{portuguese}
Exemplo em português: os anexos são documentos de terceiros, incluídos para
consulta. Para inserir um PDF, coloque-o na pasta \texttt{figures/} (ou numa
pasta \texttt{annexes/}) e descomente a linha \texttt{\textbackslash
includepdf} deste ficheiro.
\end{otherlanguage}

% PT: Inserir todas as páginas de um PDF, à escala da página.
%     pages=- : todas as páginas; pages={1-3} : só as páginas 1 a 3.
% EN: Insert every page of a PDF, scaled to the page.
%     pages=- : all pages; pages={1-3} : only pages 1 to 3.
% \includepdf[pages=-, scale=0.9, pagecommand={}]{annexes/document.pdf}
